\documentclass[11pt]{article}

\usepackage{amsmath, amssymb, amsthm}
\usepackage{geometry}
\usepackage{hyperref}
\usepackage{xcolor}
\usepackage{booktabs}
\usepackage{array}
\usepackage{parskip}
\usepackage{microtype}
\usepackage{enumitem}

\geometry{margin=1.25in}

\hypersetup{
  colorlinks=true,
  linkcolor=blue!60!black,
  citecolor=blue!60!black,
  urlcolor=blue!60!black
}

\theoremstyle{plain}
\newtheorem{definition}{Definition}[section]
\newtheorem{proposition}{Proposition}[section]
\newtheorem{remark}{Remark}[section]

\newcommand{\E}{\mathcal{E}}
\newcommand{\G}{\mathcal{G}}
\newcommand{\R}{\mathbb{R}}
\newcommand{\Crit}{\mathrm{Crit}}
\newcommand{\Min}{\mathrm{Min}}
\newcommand{\Sad}{\mathrm{Sad}}
\newcommand{\Max}{\mathrm{Max}}
\newcommand{\gCST}{\mathrm{gCST}}
\newcommand{\CST}{\mathrm{CST}}
\newcommand{\Li}{\mathrm{L}_\mathrm{i}}
\newcommand{\Lc}{\mathrm{L}_\mathrm{c}}
\newcommand{\Lj}{\mathrm{L}_\mathrm{j}}
\newcommand{\Lg}{\mathrm{L}_\mathrm{g}}
\newcommand{\Gv}{\mathrm{Gv}}
\newcommand{\BVAB}{\mathrm{BVAB1}}

\title{\textbf{Geodesic Community State Types (gCSTs):}\\
  A Morse-Stratified Framework for Vaginal Microbiome\\
  Community Classification}

\author{Pawel Gajer \and Jacques Ravel}

\date{\today}

%%%%%%%%%%%%%%%%%%%%%%%%%%%%%%%%%%%%%%%%%%%%%%%%%%%%%%%%%%%
\begin{document}
\maketitle
\tableofcontents
\newpage

%%%%%%%%%%%%%%%%%%%%%%%%%%%%%%%%%%%%%%%%%%%%%%%%%%%%%%%%%%%
\section{Motivation}
\label{sec:motivation}

Standard Community State Types (CSTs) have become the lingua franca of vaginal
microbiome research since their introduction by Ravel et al.\ (2011).  The
VALENCIA classifier~\cite{france2020} operationalised CSTs as
\emph{nearest-centroid} assignments: each of 13 reference centroids was
identified from a training corpus of 13{,}160 taxonomic profiles, and a new
sample is assigned to whichever centroid is closest under a suitable
compositional distance.

This approach captures the \emph{attractor structure} of vaginal community
space well -- it identifies the stable, species-dominated states around which
the majority of samples cluster.  However, it is silent about the
\emph{transitional structure}: the well-populated compositional paths that
connect one CST attractor to another.  Empirically, such paths are not
artefacts of small samples; they appear consistently across geographically and
ethnically diverse cohorts.  The four arms most prominently observed are:
\begin{itemize}[noitemsep]
  \item $\Li \to \Lc$: transition between \textit{L.\ iners} and
        \textit{L.\ crispatus} dominance,
  \item $\Li \to \Lj$: transition toward \textit{L.\ jensenii},
  \item $\Li \to \Gv$: transition toward \textit{G.\ vaginalis},
  \item $\Li \to \BVAB$: transition toward \textit{Ca.\ Lachnocurva vaginae}.
\end{itemize}
These arms are not adequately described by 0-dimensional centroid assignment:
samples on an arm receive low similarity scores and are often assigned
discordantly between VALENCIA and hierarchical clustering.  A principled
extension of the CST concept to 1-dimensional structures is therefore needed.

The present document develops such an extension -- \emph{Geodesic Community
State Types} (gCSTs) -- grounded in discrete Morse theory on the taxon graph.
The framework is explicitly dimensional: CSTs are the 0-dimensional layer,
gCSTs the 1-dimensional layer, and higher-dimensional generalisations are
delineated for future work.

%%%%%%%%%%%%%%%%%%%%%%%%%%%%%%%%%%%%%%%%%%%%%%%%%%%%%%%%%%%
\section{The Taxon Graph and the Evenness Field}
\label{sec:setup}

\subsection{Taxon graph}
\label{sec:graph}

Let $S = \{s_1, \ldots, s_n\}$ be a set of vaginal microbiome samples, each
represented as a composition vector $\mathbf{x}_i \in \Delta^{p-1}$ (the
$(p{-}1)$-dimensional simplex, with $p$ taxa).  We work with
$L^\infty$-normalised relative abundances throughout, but the framework is
agnostic to the specific normalisation.

\begin{definition}[Taxon graph]
  The \emph{taxon graph} $\G = (V, E, w)$ is an intersection $k$-nearest
  neighbour (ikNN) graph constructed on $S$, where
  \begin{itemize}[noitemsep]
    \item $V = S$ (samples as vertices),
    \item $(s_i, s_j) \in E$ if and only if $s_j$ is among the $k$ nearest
          neighbours of $s_i$ \emph{and} $s_i$ is among the $k$ nearest
          neighbours of $s_j$,
    \item $w_{ij}$ is a decreasing function of the compositional distance
          between $s_i$ and $s_j$.
  \end{itemize}
  The optimal $k$ is selected by minimising the generalised cross-validation
  (GCV) score of a graph spectral smoother fitted to the preterm birth outcome.
\end{definition}

\subsection{The evenness scalar field}
\label{sec:evenness}

\begin{definition}[Evenness field]
  Let $\E : V \to \R_{\geq 0}$ be a scalar field on the taxon graph defined
  vertex-wise by a measure of compositional evenness -- concretely, by the
  Shannon entropy
  \[
    \E(s_i) = -\sum_{j=1}^{p} x_{ij} \log x_{ij},
  \]
  or by any monotone transform thereof.  High values of $\E$ correspond to
  compositionally diverse, multi-taxon communities; low values correspond to
  communities strongly dominated by a single taxon.
\end{definition}

\begin{remark}[Morse convention]
  We adopt the convention that gradient flow is driven by $-\nabla \E$, i.e.\
  \emph{downhill} in evenness.  Under this convention:
  \begin{itemize}[noitemsep]
    \item \textbf{local minima} of $\E$ are the \emph{sinks} of the flow --
          they attract trajectories and correspond to single-species dominated,
          compositionally pure states;
    \item \textbf{local maxima} of $\E$ are the \emph{sources} -- they repel
          trajectories and correspond to highly diverse, no-dominant-taxon
          communities;
    \item \textbf{saddle points} are the \emph{bottlenecks} through which flow
          passes from high-evenness regions toward two distinct minima.
  \end{itemize}
  This convention is used consistently throughout the paper.
\end{remark}

%%%%%%%%%%%%%%%%%%%%%%%%%%%%%%%%%%%%%%%%%%%%%%%%%%%%%%%%%%%
\section{Discrete Morse Stratification of Community Space}
\label{sec:morse}

We work in the setting of discrete Morse theory~\cite{forman2002}, adapted to
the weighted graph $(\G, \E)$.  The \emph{critical points} of $\E$ on $\G$
are:
\begin{align*}
  \Min(\E) &= \{v \in V : \E(v) < \E(u) \text{ for all neighbours } u \sim v\},\\
  \Sad(\E) &= \{v \in V : v \text{ has exactly one ascending and one descending
             direction in } \E\},\\
  \Max(\E) &= \{v \in V : \E(v) > \E(u) \text{ for all neighbours } u \sim v\}.
\end{align*}

The \emph{Morse-Smale complex} of $(\G, \E, {-\nabla\E})$ stratifies the
vertex set $V$ into cells according to the asymptotic behaviour of gradient
flow trajectories.

\subsection{The dimensional hierarchy}
\label{sec:hierarchy}

Table~\ref{tab:hierarchy} summarises the correspondence between Morse-theoretic
objects, community space geometry, and biological interpretation.

\begin{table}[h]
  \centering
  \renewcommand{\arraystretch}{1.4}
  \begin{tabular}{clll}
    \toprule
    \textbf{Dim.} & \textbf{Critical type} & \textbf{$\E$ value} &
    \textbf{Biological interpretation} \\
    \midrule
    0 & Local minimum of $\E$ & Low  &
      Single-taxon dominated state; canonical CST core \\
    1 & Index-1 saddle of $\E$ & Intermediate &
      Transitional community; gCST arm generator \\
    2 & Local maximum of $\E$ & High &
      Diverse, no-dominant-taxon community; CST~IV-like sheet \\
    \bottomrule
  \end{tabular}
  \caption{Dimensional stratification of vaginal community space under the
    evenness field $\E$ with flow convention $-\nabla\E$.}
  \label{tab:hierarchy}
\end{table}

%%%%%%%%%%%%%%%%%%%%%%%%%%%%%%%%%%%%%%%%%%%%%%%%%%%%%%%%%%%
\section{The gCST Framework}
\label{sec:gcst}

\subsection{CST-0: canonical community state types}
\label{sec:cst0}

\begin{definition}[CST-0 core]
  A \emph{CST-0 core} is a connected component of the basin of attraction of a
  local minimum $v^* \in \Min(\E)$ under the flow $-\nabla\E$.  Formally,
  \[
    \mathrm{Basin}(v^*) =
      \{v \in V : \phi_t(v) \to v^* \text{ as } t \to \infty\},
  \]
  where $\phi_t$ denotes the gradient flow of $-\E$ on $\G$.
\end{definition}

Each CST-0 core corresponds to one of the canonical, species-dominated
community states: $\Lc$ (CST~I), $\Lg$ (CST~II), $\Li$ (CST~III),
$\mathrm{L_j}$ (CST~V), and the dominant taxa of CST~IV sub-types.

\paragraph{Relationship to VALENCIA.}
The VALENCIA centroids~\cite{france2020} are empirical summaries of CST-0
cores: each centroid is the weighted mean composition of samples assigned to
one CST by hierarchical clustering of the 13{,}160-sample training corpus.
VALENCIA's nearest-centroid classifier is therefore an approximation to
CST-0 basin assignment.  The gCST framework takes VALENCIA centroids as
\emph{valid 0-dimensional anchors} and builds higher-dimensional structure on
top of them: the graph topology of the 13{,}160-sample taxon graph defines
the 1-dimensional arms, and VALENCIA centroids serve as the fixed endpoints
(boundary conditions) of those arms.

\subsection{gCST-1: geodesic transition community state types}
\label{sec:gcst1}

\begin{definition}[gCST-1 arm]
  Let $\sigma \in \Sad(\E)$ be an index-1 saddle of the evenness field.
  The \emph{gCST-1 arm associated with $\sigma$} is the 1-dimensional unstable
  separatrix of $\sigma$ under the flow $-\nabla\E$:
  \[
    \mathrm{Arm}(\sigma) =
      W^u(\sigma) =
      \{v \in V : \phi_{-t}(v) \to \sigma \text{ as } t \to \infty\},
  \]
  i.e.\ the set of vertices whose \emph{backward} flow under $-\nabla\E$
  converges to $\sigma$.  The arm is a 1-manifold with two boundary
  points, each lying in a distinct CST-0 basin.
\end{definition}

\begin{remark}[Geodesic characterisation]
  In the discrete graph setting, $\mathrm{Arm}(\sigma)$ is operationalised as
  the \emph{shortest geodesic path} (with respect to edge weights $w$) between
  the two CST-0 endpoints linked by the saddle $\sigma$.  Each vertex along
  this path is a \emph{reference composition} for the gCST-1 arm, playing the
  role that a VALENCIA centroid plays for a CST-0.  Classification of a new
  sample to a gCST-1 arm is based on its maximum similarity to any vertex in
  the arm path:
  \[
    \mathrm{sim}(\mathbf{x},\, \mathrm{Arm}(\sigma))
      = \max_{v \in \mathrm{Arm}(\sigma)} \mathrm{sim}(\mathbf{x}, \mathbf{x}_v).
  \]
\end{remark}

\paragraph{Hard and soft classification.}
Hard gCST-1 assignment maps each sample to
$\arg\max_\sigma \mathrm{sim}(\mathbf{x}, \mathrm{Arm}(\sigma))$, subject to a
minimum similarity threshold $\tau$.  Soft (overlapping) assignment instead
assigns a membership vector
$\mathbf{m}(\mathbf{x}) \in [0,1]^{|\Sad(\E)|}$ obtained by
normalising the per-arm similarities.  Soft assignment is the preferred default
because it reflects the continuous nature of compositional transitions and is
consistent with the Morse-theoretic picture, in which the separatrix is a
continuous curve rather than a crisp boundary.

\subsection{gCST-2: community state sheets (future work)}
\label{sec:gcst2}

\begin{definition}[gCST-2 sheet -- informal]
  A \emph{gCST-2 sheet} is a 2-dimensional cell in the Morse-Smale complex of
  $(\G, \E)$ bounded by the 1-dimensional skeleton of gCST-1 arms.  It
  corresponds to a region of community space in which no single taxon dominates
  and in which samples are simultaneously drawn toward three or more competing
  CST-0 minima.  The CST~IV communities -- highly diverse, Lactobacillus-poor
  states -- are the primary biological candidates for gCST-2 characterisation.
\end{definition}

The formal development of gCST-2 classifiers, including the triangulation of
2-dimensional cells and the associated similarity scores, is deferred to future
work.

%%%%%%%%%%%%%%%%%%%%%%%%%%%%%%%%%%%%%%%%%%%%%%%%%%%%%%%%%%%
\section{The Role of \textit{L.\ iners} as Transition Hub}
\label{sec:li-hub}

A striking empirical observation motivates a formal conjecture.  The four
dominant arms identified in the PreSSMat (Zambia) and ZAPPS (Bangladesh)
cohorts all share $\Li$ as one endpoint:
\[
  \Li \leftrightarrow \Lc, \quad
  \Li \leftrightarrow \Lj, \quad
  \Li \leftrightarrow \Gv, \quad
  \Li \leftrightarrow \BVAB.
\]
This is not consistent with $\Li$ being merely one of several equivalent CST
minima: it suggests that the $\Li$ minimum has anomalously high
\emph{saddle-incidence degree} in the Morse-Smale graph.

\begin{definition}[Saddle-incidence degree]
  For a local minimum $v^* \in \Min(\E)$, define its \emph{saddle-incidence
  degree}
  \[
    \deg_\Sad(v^*) =
      \bigl|\{\sigma \in \Sad(\E) :
        v^* \in \partial\,\mathrm{Arm}(\sigma)\}\bigr|,
  \]
  i.e.\ the number of gCST-1 arms that have $v^*$ as one of their two
  boundary endpoints.
\end{definition}

\begin{proposition}[Li-hub conjecture -- testable]
  Let $\G_{13k}$ denote the taxon graph built on the VALENCIA 13{,}160-sample
  training corpus.  Then
  \[
    \deg_\Sad(\Li) > \deg_\Sad(v^*)
    \quad \text{for all } v^* \in \Min(\E) \setminus \{\Li\}.
  \]
\end{proposition}

This conjecture is biologically grounded: \textit{L.\ iners} is the species
most capable of coexisting transiently with diverse microbial partners.  It
occupies an intermediate ecological niche -- it is Lactobacillus-dominant
enough to suppress overt BV, yet metabolically less protective than
\textit{L.\ crispatus}, and sufficiently tolerant of Gardnerella, Ureaplasma,
and anaerobic species to form stable co-occurrence states.  High
saddle-incidence degree is the topological expression of this ecological
flexibility.

\paragraph{Empirical tests.}
The conjecture generates three testable predictions:
\begin{enumerate}[noitemsep]
  \item \textbf{Saddle degree}: compute $\deg_\Sad(v^*)$ for all minima in
        $\G_{13k}$; expect $\Li$ to rank first.
  \item \textbf{Arm occupancy}: among all samples assigned to gCST-1 arms,
        expect the fraction lying on $\Li$-linked arms to be highest.
  \item \textbf{Transition flux} (longitudinal data): in cohorts with
        repeated sampling, expect the transition rate through $\Li$-linked
        arms to exceed that through arms not involving $\Li$.
\end{enumerate}

%%%%%%%%%%%%%%%%%%%%%%%%%%%%%%%%%%%%%%%%%%%%%%%%%%%%%%%%%%%
\section{Building gCSTs from the VALENCIA Training Corpus}
\label{sec:implementation}

\subsection{Overview}

The construction of a reference gCST-1 classifier from the 13{,}160-sample
VALENCIA corpus proceeds in five stages.

\begin{enumerate}
  \item \textbf{Graph construction.}  Build the taxon graph $\G_{13k}$ using
        the ikNN algorithm with GCV-optimal $k$.  Apply geometric pruning to
        remove long outlier edges.

  \item \textbf{Evenness field.}  Compute $\E(v)$ for every vertex.  Identify
        $\Min(\E)$, $\Sad(\E)$, and $\Max(\E)$ using the discrete Morse
        framework.

  \item \textbf{Arm identification.}  For each saddle $\sigma \in \Sad(\E)$,
        trace the descending separatrices to their two flanking minima.
        Operationalise $\mathrm{Arm}(\sigma)$ as the geodesic shortest path
        (edge-weight metric) between those two minima, subject to the
        constraint that the path passes through $\sigma$ or its immediate
        neighbourhood.

  \item \textbf{Reference path table.}  Collect all arm-path vertices into a
        reference table analogous to the VALENCIA centroid table.  Each row is
        an observed sample composition; each row is annotated with its arm
        label (e.g.\ \texttt{arm.li.lc}, \texttt{arm.li.gv}) and its
        geodesic position along the arm (a scalar in $[0,1]$, where 0 and 1
        correspond to the two CST-0 endpoints).

  \item \textbf{Classifier.}  Given a new sample $\mathbf{x}$, compute
        $\mathrm{sim}(\mathbf{x}, \mathrm{Arm}(\sigma))$ for all arms and
        all CST-0 centroids (VALENCIA or newly derived).  Assign hard or soft
        labels as described in Section~\ref{sec:gcst1}.
\end{enumerate}

\subsection{Relationship between VALENCIA centroids and gCST-1 endpoints}

Two complementary formulations are possible and should be evaluated
empirically:

\begin{description}
  \item[Formulation A -- Graph-derived endpoints.]  CST-0 cores are defined
        entirely by the Morse minima of $\E$ on $\G_{13k}$.  VALENCIA
        centroids serve as a validation reference but not as input.

  \item[Formulation B -- VALENCIA-anchored endpoints.]  The 13 VALENCIA
        centroids are taken as fixed 0-dimensional anchors.  The gCST-1 arms
        are then defined as geodesic paths in $\G_{13k}$ whose endpoints are
        the graph vertices closest to the VALENCIA centroids.  This formulation
        ensures backward compatibility: any sample previously classifiable by
        VALENCIA retains its CST-0 label, and gCST-1 labels are assigned only
        to samples in the saddle regions.
\end{description}

Formulation~B is recommended for an initial publication because it builds
incrementally on the established VALENCIA infrastructure and makes the
extension easy to communicate to the vaginal microbiome community.

\subsection{Coverage of African cohort arms}

The VALENCIA training corpus is predominantly North American (IGS/HMP cohorts).
The $\Li \to \BVAB$ arm, which is prominent in the Zambia and Bangladesh
cohorts, may be underrepresented in $\G_{13k}$ because \textit{Ca.\
Lachnocurva vaginae} is substantially more prevalent in sub-Saharan African
women.  Two mitigations are recommended:

\begin{enumerate}[noitemsep]
  \item Verify that the $\Li \to \BVAB$ arm exists in $\G_{13k}$ (possibly
        as a short or sparse filament) before reporting its gCST-1 properties.
  \item For transferable arm definitions, construct $\G_{13k+ZB}$ by merging
        the VALENCIA corpus with the PreSSMat/ZAPPS data, define arm paths
        on the merged graph, and restrict reference profiles to the 13k
        vertices for classifier deployment.
\end{enumerate}

%%%%%%%%%%%%%%%%%%%%%%%%%%%%%%%%%%%%%%%%%%%%%%%%%%%%%%%%%%%
\section{Summary of Formal Definitions}
\label{sec:summary}

For convenience, we collect the core definitions.

\begin{description}
  \item[Scalar field.] $\E : V(\G) \to \R_{\geq 0}$, evenness (Shannon
        entropy) of sample composition.  Flow convention: $-\nabla\E$
        (downhill in evenness).

  \item[CST-0 core.] Basin of attraction of a local minimum
        $v^* \in \Min(\E)$ under $-\nabla\E$.  Corresponds to a
        single-taxon-dominated, compositionally pure community state.

  \item[gCST-1 arm.] One-dimensional unstable separatrix
        $W^u(\sigma)$ of an index-1 saddle $\sigma \in \Sad(\E)$ under
        $-\nabla\E$.  Operationalised as the geodesic shortest path between the
        two CST-0 endpoints linked by $\sigma$.  Serves as the 1-dimensional
        reference object for classification of transitional communities.

  \item[gCST-2 sheet (future).] Two-dimensional cell of the Morse-Smale
        complex, bounded by the 1-dimensional skeleton.  Corresponds to
        highly diverse, no-dominant-taxon communities (CST~IV-like states).

  \item[Saddle-incidence degree.] $\deg_\Sad(v^*)$: number of gCST-1 arms
        incident to a given CST-0 minimum.  The \textit{L.\ iners} minimum is
        conjectured to have the highest saddle-incidence degree among all
        CST-0 minima.
\end{description}

%%%%%%%%%%%%%%%%%%%%%%%%%%%%%%%%%%%%%%%%%%%%%%%%%%%%%%%%%%%
\section*{Acknowledgements}

The gCST framework emerged from ongoing analysis of the PreSSMat and ZAPPS
cohorts as part of the Gates Foundation-funded project on spontaneous preterm
birth.  We thank Michael T.\ France for the VALENCIA classifier and training
corpus, which provide the 0-dimensional foundation on which the present
extension is built.

%%%%%%%%%%%%%%%%%%%%%%%%%%%%%%%%%%%%%%%%%%%%%%%%%%%%%%%%%%%
\begin{thebibliography}{9}

\bibitem{france2020}
France~MT, Ma~B, Gajer~P, Brown~S, Humphrys~MS, Holm~JB, Waetjen~LE,
Brotman~RM, Ravel~J.
VALENCIA: a nearest centroid classification method for vaginal microbial
communities based on composition.
\textit{Microbiome} 2020;8:166.

\bibitem{ravel2011}
Ravel~J, Gajer~P, Abdo~Z, Schneider~GM, Koenig~SSK, McCulle~SL,
Karlebach~S, Gorle~R, Russell~J, Tacket~CO, Brotman~RM, Davis~CC, Ault~K,
Peralta~L, Forney~LJ.
Vaginal microbiome of reproductive-age women.
\textit{Proc Natl Acad Sci USA} 2011;108(Suppl~1):4680--4687.

\bibitem{gajer2012}
Gajer~P, Brotman~RM, Bai~G, Sakamoto~J, Schutte~UME, Zhong~X, et al.
Temporal dynamics of the human vaginal microbiota.
\textit{Sci Transl Med} 2012;4:132ra52.

\bibitem{forman2002}
Forman~R.
A user's guide to discrete Morse theory.
\textit{Sem Lothar Combin} 2002;48:B48c.

\end{thebibliography}

\end{document}
